\textbf{Phần 1. Trắc nghiệm lựa chọn}
\setcounterpageref{ex}{0}

\begin{ex} %[2D10V50]
Điền dấu thích hợp vào chỗ chấm: $700 \dots 500 + 200$.
\choice
{\True $=$}
{$>$}
{$<$}
{Không so sánh được}
\loigiai{
$500 + 200 = 700$, vậy $700 = 700$.
\textbf{Đáp án của câu hỏi là A}
}
\end{ex}

\begin{ex} %[2D10V50]
Trong các số $400, 700, 200, 900$, số lớn nhất là
\choice
{$700$}
{$400$}
{\True $900$}
{$200$}
\loigiai{
Số lớn nhất là 900.
\textbf{Đáp án của câu hỏi là C}
}
\end{ex}

\begin{ex} %[2D10V50]
Tìm số tròn trăm $x$ sao cho: $300 < x < 500$. Giá trị của $x$ là
\choice
{$350$}
{\True $400$}
{$450$}
{$600$}
\loigiai{
$x$ tròn trăm nằm giữa 300 và 500 nên $x = 400$.
\textbf{Đáp án của câu hỏi là B}
}
\end{ex}

\begin{ex} %[2D10V50]
Dãy số nào sau đây được sắp xếp theo thứ tự giảm dần?
\choice
{$200, 500, 800$}
{$600, 700, 300$}
{$800, 600, 700$}
{\True $400, 200, 100$}
\loigiai{
Dãy $400, 200, 100$ xếp giảm dần.
\textbf{Đáp án của câu hỏi là D}
}
\end{ex}

\begin{ex} %[2D10V50]
So sánh hai biểu thức: $800 - 300$ và $200 + 400$.
\choice
{\True $800 - 300 < 200 + 400$}
{$800 - 300 > 200 + 400$}
{$800 - 300 = 200 + 400$}
{Không so sánh được}
\loigiai{
$800 - 300 = 500$; $200 + 400 = 600$. Do $500 < 600$ nên $800 - 300 < 200 + 400$.
\textbf{Đáp án của câu hỏi là A}
}
\end{ex}

\begin{ex} %[2D10V50]
Cho các số tròn chục: $230, 280, 250, 210$. Số bé nhất trong dãy là
\choice
{$230$}
{$250$}
{\True $210$}
{$280$}
\loigiai{
Cùng có chữ số hàng trăm là 2, hàng chục nhỏ nhất là 1 nên 210 bé nhất.
\textbf{Đáp án của câu hỏi là C}
}
\end{ex}

\begin{ex} %[2D10V50]
Có bao nhiêu số tròn trăm $x$ thỏa mãn $200 \le x < 700$?
\choice
{$4$ số}
{\True $5$ số}
{$6$ số}
{$3$ số}
\loigiai{
Các số là: 200, 300, 400, 500, 600. Có 5 số.
\textbf{Đáp án của câu hỏi là B}
}
\end{ex}

\begin{ex} %[2D10V50]
Cây cầu thứ nhất dài $300\text{ m}$, cây cầu thứ hai dài $500\text{ m}$. Cây cầu thứ hai dài hơn cây cầu thứ nhất bao nhiêu mét?
\choice
{$100\text{ m}$}
{$300\text{ m}$}
{$400\text{ m}$}
{\True $200\text{ m}$}
\loigiai{
$500 - 300 = 200$ (m).
\textbf{Đáp án của câu hỏi là D}
}
\end{ex}

\begin{ex} %[2D10V50]
Điền số tròn trăm thích hợp vào chỗ chấm: $600 + \dots = 1 000$.
\choice
{\True $400$}
{$300$}
{$500$}
{$200$}
\loigiai{
$1 000 - 600 = 400$.
\textbf{Đáp án của câu hỏi là A}
}
\end{ex}

\begin{ex} %[2D10V50]
Một cửa hàng buổi sáng bán được $300\text{ kg}$ gạo, buổi chiều bán được nhiều hơn buổi sáng $200\text{ kg}$ gạo. Buổi chiều cửa hàng bán được bao nhiêu ki-lô-gam gạo?
\choice
{$400\text{ kg}$}
{$600\text{ kg}$}
{\True $500\text{ kg}$}
{$700\text{ kg}$}
\loigiai{
$300 + 200 = 500$ (kg).
\textbf{Đáp án của câu hỏi là C}
}
\end{ex}

\begin{ex} %[2D10V50]
Trong các khẳng định sau, khẳng định nào SAI?
\choice
{$900 > 800$}
{\True $1 000 < 900$}
{$300 + 400 = 700$}
{$500 - 100 > 300$}
\loigiai{
$1 000 > 900$, do đó $1 000 < 900$ là SAI.
\textbf{Đáp án của câu hỏi là B}
}
\end{ex}

\begin{ex} %[2D10V50]
Có ba bể nước: Bể A chứa $400\text{ l}$, bể B chứa $600\text{ l}$, bể C chứa $500\text{ l}$. Thứ tự các bể từ chứa ít nước nhất đến nhiều nước nhất là
\choice
{Bể B, Bể C, Bể A}
{Bể C, Bể A, Bể B}
{Bể A, Bể B, Bể C}
{\True Bể A, Bể C, Bể B}
\loigiai{
Vì $400 < 500 < 600$ nên thứ tự là Bể A, Bể C, Bể B.
\textbf{Đáp án của câu hỏi là D}
}
\end{ex}

\textbf{Phần 2. Trắc nghiệm đúng sai}
\setcounterpageref{ex}{0}

\begin{ex} %[2D10V50]
Kiểm tra các nhận định so sánh số tròn trăm:
\choiceTF
{1 000 là số tròn trăm lớn nhất có bốn chữ số.}
{\True Nếu a = 400 và b = 300 thì a + b < 800.}
{\True Số tròn chục liền sau số 390 là 400.}
{Hiệu giữa hai số tròn trăm liên tiếp luôn bằng 10.}
\loigiai{
\begin{itemize}
    \item Ý a) Sai: 1 000 là số tròn trăm bé nhất có bốn chữ số.
    \item Ý b) Đúng: 400 + 300 = 700 < 800. Đúng.
    \item Ý c) Đúng: 390 + 10 = 400. Đúng.
    \item Ý d) Sai: Hiệu giữa hai số tròn trăm liên tiếp luôn bằng 100.
\end{itemize}
\begin{itemize}
    \item a) \textbf{Sai}.
    \item b) \textbf{Đúng}.
    \item c) \textbf{Đúng}.
    \item d) \textbf{Sai}.
\end{itemize}
}
\end{ex}

\begin{ex} %[2D10V50]
Cho ba số tròn trăm 300, 500, 800:
\choiceTF
{\True 300 + 500 = 800.}
{\True Số lớn nhất gấp hơn 2 lần số bé nhất.}
{500 là số ở chính giữa 300 và 800.}
{\True Tổng ba số bằng 1 600.}
\loigiai{
\begin{itemize}
    \item Ý a) Đúng: Đúng.
    \item Ý b) Đúng: 800 > 300 x 2 = 600. Đúng.
    \item Ý c) Sai: Khoảng cách 500 - 300 = 200, 800 - 500 = 300 (không bằng nhau).
    \item Ý d) Đúng: 300 + 500 + 800 = 1 600. Đúng.
\end{itemize}
\begin{itemize}
    \item a) \textbf{Đúng}.
    \item b) \textbf{Đúng}.
    \item c) \textbf{Sai}.
    \item d) \textbf{Đúng}.
\end{itemize}
}
\end{ex}

\textbf{Phần 3. Trả lời ngắn}
\setcounterpageref{ex}{0}

\begin{ex} %[2D10V50]
Cho ba số tròn trăm $a, b, c$ với $a < b < c$. Biết $a$ bé nhất có ba chữ số, $c$ lớn nhất có ba chữ số và $b$ ở chính giữa. Tìm $b$.
\shortans{500}
\loigiai{
\textbf{Đáp số:} $500$\\
$a = 100, c = 900$. Số ở chính giữa là 500.
}
\end{ex}

\begin{ex} %[2D10V50]
Đoạn đường từ nhà Nam đến trường dài 700 m, từ nhà Bắc đến trường dài 400 m. Đoạn đường từ nhà Nam dài hơn từ nhà Bắc bao nhiêu mét?
\shortans{300}
\loigiai{
\textbf{Đáp số:} $300$\\
700 - 400 = 300 m.
}
\end{ex}

\begin{ex} %[2D10V50]
Tìm số tròn trăm $x$ biết $900 - x = 200 + 300$.
\shortans{400}
\loigiai{
\textbf{Đáp số:} $400$\\
$200 + 300 = 500$. $x = 900 - 500 = 400$.
}
\end{ex}

\begin{ex} %[2D10V50]
Có bao nhiêu số tròn trăm lớn hơn 150 và bé hơn 650?
\shortans{5}
\loigiai{
\textbf{Đáp số:} $5$\\
Đó là 200, 300, 400, 500, 600. Có 5 số.
}
\end{ex}

\begin{ex} %[2D10V50]
Tính hiệu của số tròn trăm lớn nhất có ba chữ số và số tròn trăm bé nhất có ba chữ số.
\shortans{800}
\loigiai{
\textbf{Đáp số:} $800$\\
900 - 100 = 800.
}
\end{ex}

\begin{ex} %[2D10V50]
Thùng thứ nhất có 300 lít dầu, thùng thứ hai có nhiều hơn 400 lít dầu. Cả hai thùng có bao nhiêu lít dầu?
\shortans{1 000}
\loigiai{
\textbf{Đáp số:} $1 000$\\
Thùng 2 có $300 + 400 = 700$ l. Cả hai thùng có $300 + 700 = 1 000$ l.
}
\end{ex}
